% ============================================================
%  J. P. Jacobsen og Vilhelm Møller, "Darwin: hans Liv og hans Lære"
%  Gyldendal, København, 1893
%  Excerpt: "Darwins Liv," pp. 38–45 (Darwins Trostab)
%  TRANSCRIPTION (Danish)
%
%  Note on authorship: The chapter "Darwins Liv" was written by
%  Vilhelm Møller, not Jacobsen. Jacobsen's own contributions to
%  the volume are the expositions of the Origin of Species and
%  (in part) the Descent of Man, which he had translated into Danish
%  in 1871–74. Møller compiled the volume posthumously and added
%  the biographical chapter and the completion of the Descent article.
%
%  Note on text: Transcribed from the Internet Archive OCR of the 1893
%  edition (darwinhanslivafj00jacouoft). OCR errors corrected by
%  contextual reconstruction; orthography follows the original.
% ============================================================
\documentclass[12pt,a4paper]{article}
\usepackage[utf8]{inputenc}
\usepackage[T1]{fontenc}
\usepackage[danish]{babel}
\usepackage{microtype}
\usepackage{setspace}
\usepackage[margin=1.2in]{geometry}
\usepackage{fancyhdr}
\usepackage{hyperref}

\hypersetup{
  pdftitle={Darwin: hans Liv og hans Lære (uddrag)},
  pdfauthor={J. P. Jacobsen og Vilhelm Møller},
  hidelinks
}

\pagestyle{fancy}
\fancyhf{}
\rhead{Møller, \emph{Darwins Liv} (1893), s.\ 38--45}
\cfoot{\thepage}

\setstretch{1.3}

\title{\textbf{Darwin: hans Liv og hans Lære}\\[0.4em]
\large Uddrag: »Darwins Trostab«\\ (»Darwins Liv,« s.\ 38--45)}
\author{J.\ P.\ Jacobsen og Vilhelm Møller}
\date{Gyldendal, København, 1893}

\begin{document}
\maketitle
\thispagestyle{fancy}

\bigskip

\noindent Darwin tror ikke, at den religiøse Følelse nogetsteds har
været stærkt udviklet hos ham. Dog var den ikke sjælden og
kunde blive levende: han siger selv, at der var Tider, da der kom
Liv i den. »Den Sindstilstand,« siger han, »som storartede Scener
vakte hos mig, og som var meget inderlig forbunden med en Tro
paa Gud, afveg ikke væsentlig fra dét, der nu om Stunder kaldes
Følelsen for det Ophøjede«; og han sammenstiller Sindstilstanden
med »de mægtige, i Aarsag ubestemte, men lignende Stemninger,
som Musik vækker«. Med andre Ord, hans religiøse Forestillinger
var associerede med, de bares oppe af en særegen Art »Lyrik«,
der var i ham: af hans Begejstring, hans Evne til at betages,
gribes, glide studsende, beundrende hen i Hengivelse.\footnote{Se
G.\ Brandes, \emph{Naturalismen i England}, s.\ 60~ff.}

»Nu derimod« — tilføjer han saa, og dette »nu« turde referere
sig allerede til den sidste Del af Rejsetiden — »nu derimod
vilde selv de mest storartede Scener ikke lade den Slags
Overbevisninger og Stemninger opstaa i min Sjæl«. Nej, ti
»storartede Scener« eksisterede overhovedet ikke længer for ham;
han fornam dem ikke længer som storartede. En eneste Gang endnu
(i Maj 1838) synes han under en Kirkekoncert, at et Kor »fik
virkelig Murene til at ryste«; ellers vil man i alle hans Breve
fra efter Rejsen ikke finde bevaret noget helt ud »ophøjet«
Indtryk, kun idylliske. En vis Drivkraft i ham var som sagt
bleven slappet. Før havde den, under et mægtigt Pres, drevet
hans Stemning som en regnbueglitrende Springvandsstraale højt,
højt op i Æterens Blaa: og paa de øverste Perler havde sitret i
Solen et gyldent Glimtende, — en Stjærne, — Stjærnen, som ledte
vise Mænd til Østerland. Nu mægtede denne Drivkraft, naar den
anspændtes, kun at sende Stemningen lavt til Vejrs, og hvad der
havde syntes en himmelsk Stjærne saas nu som et Guldæble, skabt
af Mennesketanker!

Men dét, Mennesker har skabt, ræsonnerer Menneskene over. —
»Om Bord paa \emph{Beagle} havde jeg været helt ortodoks \ldots\
og citeret Biblen som uigendrivelig Autoritet i moralske
Spørgsmaal \ldots\ Men paa denne Tid, det vil sige 1836--39,
var jeg efterhaanden kommen til at indse, at der ikke var mere
Grund til at tro det gamle Testamente end Hinduernes hellige
Skrifter \ldots\ Ved fortsat Eftertanke tvang Omstændighederne
mig lidt efter lidt til ikke mer at tro paa Kristendommen som en
guddommelig Aabenbaring.«

I hans Bevidsthed udpønsedes mulige Bevismidler for Bibelens
Sandhed; men saa snart Kristendommen ikke mere stod uigendrivelig
fast, formaaede disse kunstige Hjælpemidler ikke at holde ved
Overbevisningen. »Ikke-Troen sneg sig saare langsomt ind«,
skriver han, »og var tilsidst fuldstændig. Det kom ikke over mig
med Stormvejrs Hurtighed eller med noget særligt Voldsomt.«

Hermed fjærnedes tillige en Opfattelse, som — ikke nødvendigt,
men vanemæssigt — var og er forbunden med den gængse kristne
Gudstro. Det er Opfattelsen af Gud som direkte Frembringer af
alle det organiske Livs Former. I et Brev fra en noget senere
Tid (1859) skriver Darwin med Føje, at ligesom Forestillingen
om, at Torden og Lyn er Følger af sekundære Aarsager, har mange
beklaget at skulle opgive, saaledes beklager mange nu at skulle
opgive Forestillingen om Gud som den umiddelbare Ophav til
de organiske Former. For de kosmiske Fænomener og for det
anorganiskes Fænomener havde Flertallet af det 19de Aarhundredes
Kristne dog ophørt at tro, at Gud var den umiddelbare Aarsag;
men med Hensyn til det organiske Livs Fænomener troede og tror
de, at Gud er det. Den gængse kristne Gudstro var og er
vanemæssigt forbunden med Tro paa den organiske Naturs
guddommelige »Planmæssighed«: og med den kan ingen gennemført
Evolutionshypotese forliges. Først altsaa da Darwin var kommen
bort — til Stadighed ikke just fra Troen paa en Gud, men for
bestandig fra det gængse Gudsbegreb — først da var det muligt
for ham at tænke sig Arterne dannede ved den gennem
Tilværelseskampen realiserede Udvælgelse af tilfældige
Variationer. Ti saavel med Ordene »Tilværelseskamp« og
»Kvalitetsvalg« (naturlig Udvælgelse) som med Ordet »tilfældig«
skydes en lang naturlig Aarsagsrække ind mellem Fænomenerne og
Gud, — med Ordet »Tilfælde« jo endog en saa lang, at
Menneskeøjet ikke mægter at forfølge mer end dens alleryderste
Led\ldots

Maaske ét Forhold endnu bør drages frem til Forklaring af
»Darwinismens« Tilblivelse. — Der er ikke i Darwins Breve og
Optegnelser nogen ligefrem Tilkendegivelse om, af hvad Stemning
hans rationalistiske Kritik blev født. Men saa ofte han tvivlede
om en Guds Tilværelse (og længer end til Tvivl gik han aldrig,
end ikke naar han gik yderligst!) — bestandig støttede han sig
til denne, og kun til denne Indvending: »der er saa megen Lidelse
i Verden!« Dét maa have været det oprindeligste og dybeste
Vantros-Motiv i hans Bevidsthed.

Hans Medfølelse med alt lidende gav ligeledes til dette Vaaben
i Kampen med Forsynstanken. Hans Sygelighed og hans til Tider
arrige Lede ved Forholdenes Uretfærdighed skærpede Vaabnet.
Det maa, selvfølgelig, have virket bedst paa Tider, hvor han
»saa sort« paa Livet, da han var i en Stemning, saa Lidelsen
her i Verden syntes ham det overvejende, Glæden det
forsvindende. Imidlertid — dersom denne Stemning havde sat sig
fast, hvis han altsaa var bleven »Pessimist«, saa kunde han
umuligt have tilegnet sig og skabt den ledende Idé, han fik ved
Læsningen af Malthus' Bog. Ti for at én skal tilegne sig og
skabe ledende Ideer af dét, han læser eller møder paa sin Vej,
er det jo ikke nok, at Forskertrangen er stærk hos ham: der maa
tillige, mellem det han sidder med og det han optager, være hin
hemmelighedsfulde Paarørenhed, som Kemien med Hensyn til
Stofferne kalder »kemisk Slægtskab«. Og den Idé, han fik, var
en lykkelig, en gjort optimistisk Idé.

Han havde omgaaedes meget med Opdrættere og hørt om deres
Gerning: nu opfattede han, kortelig sagt, Tilværelseskampen som
den vældige og stilfærdige »Opdrætter« af det organiske Livs
Former! Men allerede Forestillingen »Opdrætning« har ligesom
festlig Glans omkring sig. Den minder, det er sandt, om Slid og
glippede Haab og fejlslagne Forsøg og mange Slags Ofre; men
langt, langt stærkere lyser den dog ud over Iagttagelsens og
Eksperimentets Glæde og peger mod det lykkeligt naaede Maal:
at der dannes en ny Varietet, nyttig og tilfredsstillende for
den, som dannede den! I Virkeligheden: bag det Darwinske
Begreb »Kvalitetsvalg« (naturlig Udvælgelse, naturlig
Opdrætning) dæmrer en Anelse om Noget eller Nogen —
»Naturen«, »Gud«, et Ufatteligt og Navnløst! — der glæder
sig over, at Livets Former ændres og fuldkommengøres til at
passe efter de skiftende Forhold. Der er forbundet med hint
Begreb en Tilfredshed over et uendeligt og universelt
Nytte-Fremskridt\ldots\ Og nu for Individernes Vedkommende,
for hver enkelt i det Livets Mylder, som gennemgaar
Tilværelseskampens Opdrætter-Proces? Ja, ogsaa for
Individerne — eller for det langt overvejende Flertal af dem —
forudsætter Darwins »Opdræt\-nings«-Idé et Overskud af Lykke.
Ti den Opdrætten maatte hurtigt tage Ende, hvor de opdrættede
var mere syge end sunde, havde mere Lede for end Lyst til
Næring, var mere uvillige end villige til at forplante sig
osv.\ osv.

Man har da ogsaa set, at just i det Halvaar, da hans
Hovedteori blev til, befandt Darwin selv sig i en ligesom dyb
og sikker Glæde. Og det er let at gætte sig til Grunden
dertil. I Oktober 1838 var det, at han læste Malthus' Bog,
— — den 29.\ Januar 1839 holdtes hans Bryllup med sin Kusine,
den udmærkede Datter af den udmærkede Josiah Wedgwood, Emma
Wedgwood! Saa vidt har da Gud Amor haft en Finger med i
»Arternes Oprindelse«. Uden Darwins lykkelige Kærlighed vilde
»Darwins Teori« ikke være bleven til.

\end{document}
